\documentclass[12pt,a4paper]{article}
\usepackage[utf8]{inputenc}
\usepackage[spanish]{babel}
\usepackage{graphicx}
\usepackage{amsmath}
\usepackage{hyperref}
\usepackage{listings}
\usepackage{xcolor}
\usepackage{geometry}
\usepackage{float}

\geometry{margin=2.5cm}

% Configuración de listings para fragmentos de código
\lstset{
    basicstyle=\ttfamily\small,
    keywordstyle=\color{blue}\bfseries,
    commentstyle=\color{green!60!black},
    stringstyle=\color{red},
    showstringspaces=false,
    breaklines=true,
    frame=single,
    numbers=left,
    numberstyle=\tiny\color{gray}
}

\title{
    \textbf{GSEA - Gestión Segura y Eficiente de Archivos}\\
    \large Documento Técnico\\
    \vspace{0.5cm}
    \large Examen Final Práctico - Sistemas Operativos
}

\author{
    Juan Pablo Rúa Cartagena\\
    Juan José Restrepo Higuita\\
    Fredy Cadavid Franco\\[0.3cm]
    Universidad EAFIT\\
    Escuela de Ciencias Aplicadas e Ingeniería
}
\date{Noviembre 2025}

\begin{document}

\maketitle
\newpage

\tableofcontents
\newpage

\section{Introducción}

GSEA (Gestión Segura y Eficiente de Archivos) es una utilidad de línea de comandos desarrollada en C++ para comprimir/descomprimir y encriptar/desencriptar archivos de manera sencilla. Se implementan algoritmos propios de compresión (Huffman) y cifrado simétrico (AES-128 en modo CBC) y se utilizan llamadas directas al sistema operativo para la entrada/salida de archivos.

\section{Objetivos}
\begin{itemize}
    \item Implementar compresión/descompresión sin pérdida.
    \item Implementar cifrado/descifrado simétrico básico.
    \item Usar llamadas directas al sistema operativo para gestionar archivos.
    \item Procesar múltiples archivos en paralelo aprovechando concurrencia a nivel de hilo.
\end{itemize}

\section{Tecnologías Utilizadas}
\begin{itemize}
    \item \textbf{Lenguaje:} C++17.
    \item \textbf{Plataformas:} Windows (API Win32) y POSIX (Linux, macOS).
    \item \textbf{Concurrencia:} Windows Threads y pthreads.
    \item \textbf{Compilación:} Make / mingw32-make, GCC/Clang/MSVC.
\end{itemize}

\section{Diseño de la Solución}

\subsection{Arquitectura general}
El sistema se organiza en módulos independientes:
\begin{itemize}
    \item \textbf{main.cpp:} Orquestación y rutas de entrada/salida.
    \item \textbf{ArgParser:} Parsing y validación de argumentos.
    \item \textbf{FileManager:} I/O mediante llamadas del SO.
    \item \textbf{Worker:} Procesamiento concurrente (un hilo por archivo de primer nivel).
    \item \textbf{HuffmanCompressor:} Compresión/descompresión.
    \item \textbf{AES:} Cifrado/descifrado AES-128 en modo CBC.
\end{itemize}

\subsection{Flujo de datos}
\begin{enumerate}
    \item El usuario invoca el programa con las opciones deseadas.
    \item ArgParser valida operaciones, rutas y clave (si aplica).
    \item main.cpp determina si la entrada es archivo o directorio.
    \item Para directorios, se listan solo los archivos de primer nivel (no recursivo) y se asigna un hilo a cada uno.
    \item Cada hilo lee el archivo, aplica operaciones en orden y escribe la salida.
    \item El hilo principal espera a que todos los hilos terminen y reporta el resultado.
\end{enumerate}

\section{Módulos}

\subsection{FileManager (file\_manager.cpp)}
Realiza la I/O con llamadas nativas:
\begin{itemize}
    \item \textbf{Windows:} CreateFileA, ReadFile, WriteFile, CloseHandle, FindFirstFileA/FindNextFileA.
    \item \textbf{POSIX:} open, read, write, close, opendir/readdir, mkdir.
\end{itemize}
Se añadió validación de tamaños con GetFileSizeEx (Windows) y verificación de \texttt{errno} en POSIX.

\subsection{Worker (worker.cpp)}
Procesa archivos en paralelo. Crea un hilo por archivo de primer nivel y consolida resultados. No implementa pool de hilos, por lo que directorios con miles de archivos pueden agotar recursos del sistema.

\subsection{HuffmanCompressor (huffman.cpp)}
Implementa Huffman con min-heap. Se añadió una validación extra para detectar datos comprimidos truncados (comprobación de bits esperados vs. disponibles).

\subsection{AES (aes.cpp)}
Cifrado/descifrado AES-128 en modo CBC. Validaciones añadidas:
\begin{itemize}
    \item Rechazo de tamaños de payload que no sean múltiplos del bloque al descifrar.
    \item Control de mínimo tamaño (requiere IV + al menos un bloque).
\end{itemize}
\textbf{Limitaciones actuales:} derivación de clave casera (XOR repetido), IV pseudoaleatorio con \texttt{std::random\_device}/\texttt{mt19937} y sin autenticación de integridad (no hay MAC/AEAD). No debe usarse para seguridad fuerte; se recomienda PBKDF2/scrypt + HMAC o un modo AEAD (GCM/ChaCha20-Poly1305) como trabajo futuro.

\section{Justificación de algoritmos}

\subsection{Compresión: Huffman}
Se eligió por ser sin pérdida, bien documentado y eficiente para datos con redundancia. Complejidad: $O(n + k \log k)$ para construir y codificar, con $k$ símbolos únicos. El formato incluye \textit{magic bytes}, tabla de frecuencias y cuenta total de bits, permitiendo reconstrucción del árbol.

\subsection{Cifrado: AES-128}
Estándar reconocido y eficiente. En este proyecto se implementa modo CBC con padding PKCS\#7. Debido a la derivación de clave simplificada y falta de autenticación, el nivel de seguridad efectivo es básico y adecuado solo para fines académicos.

\section{Implementación y consideraciones}
\begin{itemize}
    \item El manejo de extensiones ahora elimina correctamente \texttt{.aes} y \texttt{.huf} cuando se usa \texttt{-ud} sobre archivos \texttt{.huf.aes}.
    \item Se valida la longitud de datos cifrados y se detectan bitstreams Huffman truncados.
    \item La lectura/escritura carga cada archivo completo en memoria; para archivos muy grandes se podría streamizar en trabajo futuro.
    \item El listado de directorios no es recursivo y el nombre de salida se aplana (solo nombre de archivo, sin subdirectorios).
\end{itemize}

\section{Estrategia de concurrencia}
\begin{itemize}
    \item Paralelismo a nivel de archivo (un hilo por archivo listado).
    \item Windows: CreateThread + WaitForMultipleObjects; POSIX: pthread\_create + pthread\_join.
    \item No hay límite de hilos; recomendable añadir pool para directorios masivos.
\end{itemize}

\section{Guía de uso}

\subsection{Requisitos}
\begin{itemize}
    \item \textbf{Windows:} MinGW-w64 o MSVC con C++17.
    \item \textbf{Linux/macOS:} GCC 7+ o Clang 5+ con pthreads.
    \item Memoria: suficiente para cargar cada archivo a procesar.
\end{itemize}

\subsection{Compilación}
\begin{lstlisting}[language=bash]
mingw32-make   # Windows (MinGW)
make           # Linux/macOS
\end{lstlisting}
El ejecutable queda en \texttt{bin/gsea.exe} o \texttt{bin/gsea}.

\subsection{Ejemplos}
\begin{lstlisting}[language=bash]
bin/gsea -c  -i datos.txt        -o datos.huf
bin/gsea -d  -i datos.huf        -o datos.txt
bin/gsea -e  -i secreto.pdf      -o secreto.aes -k clave123
bin/gsea -ce -i carpeta/         -o salida/     -k clave123   # no recursivo
bin/gsea -ud -i salida/          -o restaurado/ -k clave123   # elimina .aes/.huf
\end{lstlisting}

\section{Pruebas}
Se incluyen scripts de prueba (\texttt{test\_build.sh} y \texttt{test\_build.bat}) que limpian, compilan y ejecutan un flujo básico de \texttt{-c/-d} y \texttt{-e/-u} verificando integridad con un archivo de texto repetitivo.

\section{Trabajo futuro}
\begin{itemize}
    \item Pool de hilos y procesamiento recursivo de directorios.
    \item Stream de archivos grandes para reducir uso de memoria.
    \item Derivación de clave robusta (PBKDF2/scrypt) e inclusión de MAC o modo AEAD.
    \item Progreso y manejo más férreo de errores E/S.
\end{itemize}

\section{Referencias}
\begin{enumerate}
    \item Huffman, D.A. (1952). ``A Method for the Construction of Minimum-Redundancy Codes''.
    \item NIST (2001). ``FIPS 197: Advanced Encryption Standard (AES)''.
    \item Stevens, W.R., \& Rago, S.A. (2013). \textit{Advanced Programming in the UNIX Environment}.
    \item Russinovich, M.E., Solomon, D.A., \& Ionescu, A. (2012). \textit{Windows Internals}.
\end{enumerate}

\end{document}
